\scanpage{116}

c'est la partie ``universellement géom'' de $\mathcal E_L^{\overline L}$, \deleted{noté $\mathcal E_{L/K}^{\unsure{g},\overline L}$} \deleted{qui \unsure{couvre} les \unsure{hom|bons}} \inserted{(relativement à $K$),} qui classifie les schémas (finis) étales sur $L$ qui sont ``non ramifiés'' sur \uline{tout} modèle \deleted{\unsure{étal}} régulier (propre ou non) de $L/K$.

\margin{noté $\mathcal E_L^{\unsure{g},\overline L}$ ou $\mathcal E_L^{\unsure{g(K)},\overline L}$}

Si $U$ est un modèle quelconque; il se plonge dans un $X$, et on a des hom surjectifs
\deleted{\unsure{q}}
\[
(21)\qquad \mathcal E_L^{\overline L}\longrightarrow\mathcal E_U^{\overline L}\longrightarrow\mathcal E_X^{\overline L}=\mathcal E_L^{\unsure{g},\overline L}
\]

On aura \deleted{\unsure{alors}}
\[
(22)\qquad U=X\setminus Z
\]
$Z$ partie ferm. de $X$ --- soient $D_i$ ses composantes de codim $1$, $x_i$ les pts gén. de $D_i$. On a des groupes de décomposition et d'inertie $N_{x_i'}$ et $I_{x_i'}$ \inserted{de $\mathcal E_U^{\overline L}$}, (relatifs au choix d'un $x_i'$ au dessus de $x_i$ dans le normalisé $X'$ dans $\overline L$), et on trouve

\margin{115}
